\documentclass[10pt]{article}
\usepackage[margin=.8in]{geometry}\usepackage[T1]{fontenc}
\usepackage{lmodern,booktabs,longtable,amsmath,hyperref}
\title{Supplement S7: Question selection and answer-budget sensitivity}
\author{Anonymous authors}\date{14 November 2026}
\begin{document}\maketitle
\section{Scope, methods and provenance}
These exploratory analyses reuse the same 600 previously examined E3 participants
from IPIP and 8,320 fitting people. Each fitting person supplies all 60 target
responses for supervised question selection. The answer budget counts observed
history answers from an evaluation person, not total training supervision.
Order construction and decoder tuning are repeated within training folds;
E3 is excluded from these fits but has influenced the research direction.
This is retrospective masking of complete questionnaires, not prospective
administration or independent confirmation.

Scalar target selection centers history inputs and forms the 240 cumulative
target indicators. With $G=X^TX+100I$ and $C=X^TT$, it greedily maximizes
$\|C_i\|^2/G_{ii}$ and updates residual quantities by a Schur complement.
The selection penalty is distinct from the decoder penalty, retuned at each
budget. Strong comparisons use the historical prior/ridge/kNN/ExtraTrees library;
target-directed forms select ridge. Historical raw-answer ridge already uses
one-hot inputs and isotonic CDF projection. The separate encoding ablation uses
common sorted/clipped CDFs. Its gain must not be added to the strong comparison.

In the fixed-form contrast labels, ``raw'' denotes the scalar-encoding ablation
and ``category'' denotes one-hot encoding; ``strong'' denotes the historical
candidate library, whose raw-answer ridge already uses one-hot encoding.
These labels must not be confused with the scalar representation used to learn
the question order in the main strong comparison.

The 32 random-form tasks 28416\_1--32 completed; they are forms on common people,
not independent replications. Grouped-selector job 28444 completed after job28442
failed on a singular selected block; the correction excludes already selected
blocks from the solve. The failed engineering receipt is retained. Curve job28457
completed. Exported grouped/curve aggregates were byte-matched to Euler results.
This manuscript update performs no scientific rerun.

\section{Fixed-form and mechanism contrasts}
Differences are first method minus reference; negative favors the first method.
Intervals below are pointwise, paired by participant and conditional on fixed
fits and question sets. The original within-experiment Holm tests address zero
differences, not a multiplicity-controlled practical-benefit claim. Full-precision
estimates, p-values and original flags remain in the accompanying JSON.
At five answers, target versus coverage has an interval whose upper bound is
above $-0.001$; no useful-gain claim at that margin follows. Grouped and adaptive
null/mixed outcomes are retained, without equivalence claims.
\subsection{Fixed forms and encoding}
\begin{longtable}{p{.49\linewidth}rrr}\toprule Contrast & $\Delta$ RPS & CI low & CI high\\\midrule\endhead
m5/target\_vs\_random\_raw & -0.00594 & -0.00711 & -0.00474\\
m5/target\_vs\_coverage\_raw & -0.00167 & -0.00268 & -0.00074\\
m5/category\_vs\_raw\_random & -0.00118 & -0.00152 & -0.00087\\
m5/category\_vs\_raw\_target & -0.00157 & -0.00205 & -0.00109\\
m5/target\_strong\_vs\_random\_strong & -0.00632 & -0.00766 & -0.00498\\
m20/target\_vs\_random\_raw & -0.00656 & -0.00734 & -0.00579\\
m20/target\_vs\_coverage\_raw & -0.00212 & -0.00269 & -0.00156\\
m20/category\_vs\_raw\_random & -0.00265 & -0.00309 & -0.00222\\
m20/category\_vs\_raw\_target & -0.00357 & -0.00410 & -0.00305\\
m20/target\_strong\_vs\_random\_strong & -0.00750 & -0.00839 & -0.00665\\
\bottomrule\end{longtable}
\subsection{Grouped versus scalar acquisition}
\begin{longtable}{p{.49\linewidth}rrr}\toprule Contrast & $\Delta$ RPS & CI low & CI high\\\midrule\endhead
m5/grouped\_vs\_scalar & 0.00026 & -0.00077 & 0.00131\\
m5/grouped\_vs\_marginal & -0.00227 & -0.00314 & -0.00138\\
m5/scalar\_vs\_marginal & -0.00253 & -0.00349 & -0.00158\\
m20/grouped\_vs\_scalar & 0.00022 & -0.00020 & 0.00065\\
m20/grouped\_vs\_marginal & -0.00537 & -0.00623 & -0.00451\\
m20/scalar\_vs\_marginal & -0.00559 & -0.00648 & -0.00471\\
\bottomrule\end{longtable}
\subsection{Adaptive selection controls}
\begin{longtable}{p{.49\linewidth}rrr}\toprule Contrast & $\Delta$ RPS & CI low & CI high\\\midrule\endhead
m5/adaptive\_vs\_fixed & 0.00055 & -0.00021 & 0.00131\\
m5/adaptive\_vs\_random\_same\_decoder & -0.00480 & -0.00579 & -0.00381\\
m5/adaptive\_vs\_archived\_strong & -0.00206 & -0.00309 & -0.00103\\
m5/adaptive\_vs\_fixed\_target\_strong & 0.00426 & 0.00315 & 0.00537\\
m20/adaptive\_vs\_fixed & -0.00003 & -0.00059 & 0.00055\\
m20/adaptive\_vs\_random\_same\_decoder & -0.00155 & -0.00219 & -0.00089\\
m20/adaptive\_vs\_archived\_strong & 0.00741 & 0.00650 & 0.00832\\
m20/adaptive\_vs\_fixed\_target\_strong & 0.01492 & 0.01373 & 0.01612\\
\bottomrule\end{longtable}
\section{Thirty-two additional random forms}
Each random form retunes the same candidate library on FIT. No form has a smaller
E3 point loss than the target-directed form at either budget. This rank statement
is not a p-value or evidence from 32 independent samples. Target RPS is
0.13410 at five answers and 0.11699 at twenty.
All per-form candidate selections are retained in random-forms.json.
\begin{longtable}{rrrr}\toprule Seed & RPS ($m=5$) & RPS ($m=20$) & Selected family (5/20)\\\midrule\endhead
1 & 0.13827 & 0.12174 & ridge/ridge\\
2 & 0.14100 & 0.12302 & ridge/ridge\\
3 & 0.14089 & 0.12393 & ridge/ridge\\
4 & 0.13766 & 0.12325 & ridge/ridge\\
5 & 0.14031 & 0.12195 & ridge/ridge\\
6 & 0.13975 & 0.12543 & ridge/ridge\\
7 & 0.13983 & 0.12350 & ridge/ridge\\
8 & 0.14018 & 0.12258 & ridge/ridge\\
9 & 0.13918 & 0.12299 & ridge/ridge\\
10 & 0.13973 & 0.12505 & ridge/ridge\\
11 & 0.13950 & 0.12239 & ridge/ridge\\
12 & 0.14097 & 0.12404 & ridge/ridge\\
13 & 0.14207 & 0.12516 & ridge/ridge\\
14 & 0.14030 & 0.12298 & ridge/ridge\\
15 & 0.14186 & 0.12577 & ridge/ridge\\
16 & 0.13793 & 0.12257 & ridge/ridge\\
17 & 0.13969 & 0.12509 & ridge/ridge\\
18 & 0.13764 & 0.12299 & ridge/ridge\\
19 & 0.13886 & 0.12158 & ridge/ridge\\
20 & 0.14170 & 0.12362 & ridge/ridge\\
21 & 0.14122 & 0.12428 & ridge/ridge\\
22 & 0.13891 & 0.12159 & ridge/ridge\\
23 & 0.13964 & 0.12166 & ridge/ridge\\
24 & 0.14252 & 0.12516 & ridge/ridge\\
25 & 0.13904 & 0.12439 & ridge/ridge\\
26 & 0.14056 & 0.12527 & ridge/ridge\\
27 & 0.14281 & 0.12556 & ridge/ridge\\
28 & 0.13857 & 0.12624 & ridge/ridge\\
29 & 0.13906 & 0.12323 & ridge/ridge\\
30 & 0.14300 & 0.12498 & ridge/ridge\\
31 & 0.13919 & 0.12378 & ridge/ridge\\
32 & 0.14013 & 0.12290 & ridge/ridge\\
\bottomrule\end{longtable}
\section{Curve against the archived 60-answer Granite hybrid}
All twenty examined budgets are reported. Positive differences favor the archived
hybrid, whose mean RPS is 0.13153. The first negative
mean difference is at seven answers, but its interval includes zero. The first
interval entirely below zero is at eight; the first entirely below $-0.001$ is
at nine. These are different criteria, not interchangeable crossing estimates.
The historical random-order interpolation of 11.44 remains a separate result.

The stored curve uses 30,000 shared paired participant resamples with seed 20261115
and percentile endpoints 0.00125/0.99875 (Bonferroni across 20 cut points).
These are approximate, conditional adjusted bootstrap intervals, not a guarantee
of full-procedure coverage. The seed is an identifier, not an execution date.
Fits, learned question order and historical hybrid recipes remain fixed.
\begin{longtable}{rrrrr}\toprule Answers & Statistical RPS & $\Delta$ vs hybrid60 & CI low & CI high\\\midrule\endhead
1 & 0.14668 & 0.01514 & 0.01348 & 0.01694\\
2 & 0.14232 & 0.01078 & 0.00905 & 0.01253\\
3 & 0.13920 & 0.00766 & 0.00595 & 0.00942\\
4 & 0.13682 & 0.00529 & 0.00359 & 0.00702\\
5 & 0.13410 & 0.00257 & 0.00079 & 0.00437\\
6 & 0.13254 & 0.00100 & -0.00080 & 0.00283\\
7 & 0.13089 & -0.00064 & -0.00251 & 0.00126\\
8 & 0.12918 & -0.00235 & -0.00419 & -0.00043\\
9 & 0.12741 & -0.00413 & -0.00595 & -0.00229\\
10 & 0.12613 & -0.00540 & -0.00725 & -0.00360\\
11 & 0.12496 & -0.00657 & -0.00841 & -0.00476\\
12 & 0.12400 & -0.00753 & -0.00945 & -0.00567\\
13 & 0.12299 & -0.00854 & -0.01045 & -0.00667\\
14 & 0.12153 & -0.01000 & -0.01191 & -0.00812\\
15 & 0.12058 & -0.01095 & -0.01286 & -0.00913\\
16 & 0.11963 & -0.01190 & -0.01376 & -0.01013\\
17 & 0.11906 & -0.01247 & -0.01433 & -0.01068\\
18 & 0.11821 & -0.01332 & -0.01515 & -0.01152\\
19 & 0.11751 & -0.01402 & -0.01586 & -0.01220\\
20 & 0.11699 & -0.01454 & -0.01635 & -0.01273\\
\bottomrule\end{longtable}
\section{Boundaries of the comparison}
The target-selected statistical predictor at five answers is worse than the
60-answer archived hybrid. At twenty it is better. Response budgets and total
historical supervision differ: the statistical selector uses all target labels
of 8,320 fitting people, while the hybrid has its own search, prior, selection
and calibration history. These results do not measure LLM value at matched five
or twenty answers or establish saved survey administrations. There are no new
small-budget LLM forecasts. The separate equal60 audit remains unchanged.

The fitting-budget and facet-split diagnostics are retained as aggregates
(fit-result.json and splits-result.json); their common participant cohort must
not be counted as independent replications. No result here changes historical
R2 Table1, its14 primary contrasts/686 endpoint rows, original Holm decisions,
instruction comparisons or CAL-selection sensitivity. No new external panel
is admitted and the A18 search stop is unchanged.
\end{document}
